\documentclass[12pt]{article}

\usepackage[utf8]{inputenc}
\usepackage[margin=1in,headheight=16pt]{geometry}
\usepackage{amsmath}
\usepackage{amssymb}
\usepackage{amsthm}
\usepackage[english]{babel}
\usepackage{fancyhdr}
\usepackage{titling}
\usepackage{needspace}
\usepackage[colorlinks=true,allcolors=blue]{hyperref}

\pagestyle{fancy}
\fancyhf{}
\rhead{The Minimal Leaf of $T^*(G/U)^{\mathrm{aff}}$}
\cfoot{\thepage}

\newtheoremstyle{paperplain}
    {6pt plus 1pt minus 1pt}{6pt plus 1pt minus 1pt}
    {\itshape}{}{\bfseries}{.}{.5em}{}
\newtheoremstyle{paperdefinition}
    {6pt plus 1pt minus 1pt}{6pt plus 1pt minus 1pt}
    {\normalfont}{}{\bfseries}{.}{.5em}{}
\theoremstyle{paperplain}
\newtheorem{theorem}{Theorem}[section]
\newtheorem{lemma}[theorem]{Lemma}
\newtheorem{proposition}[theorem]{Proposition}
\newtheorem{corollary}[theorem]{Corollary}
\theoremstyle{paperdefinition}
\newtheorem{definition}[theorem]{Definition}
\newtheorem{remark}[theorem]{Remark}
\AddToHook{env/proposition/before}{\Needspace{5\baselineskip}}
\AddToHook{env/proposition/begin}{\samepage}
\AddToHook{env/theorem/before}{\Needspace{6\baselineskip}}
\AddToHook{env/theorem/begin}{\samepage}
\AddToHook{env/lemma/before}{\Needspace{5\baselineskip}}
\AddToHook{env/lemma/begin}{\samepage}
\AddToHook{env/definition/before}{\Needspace{5\baselineskip}}
\AddToHook{env/definition/begin}{\samepage}

\DeclareMathOperator{\Spec}{Spec}
\DeclareMathOperator{\Sym}{Sym}
\DeclareMathOperator{\Ad}{Ad}
\newcommand{\C}{\mathbb C}
\newcommand{\Omin}{\mathcal O_{\min}}
\newcommand{\g}{\mathfrak g}
\newcommand{\tfrak}{\mathfrak t}

\hypersetup{
    pdftitle={The Minimal Leaf of T*(G/U)\textasciicircum aff},
    pdfauthor={GPT-6 Astra},
    pdfsubject={A Poisson embedding of the minimal nilpotent orbit closure}
}

\title{The Minimal Leaf of \texorpdfstring{$T^*(G/U)^{\mathrm{aff}}$}{T*(G/U)\textasciicircum aff}}
\author{GPT-6 Astra}
\date{}

\setlength{\droptitle}{-4.25em}
\pretitle{\centering\LARGE}
\posttitle{\par}
\preauthor{\centering\large}
\postauthor{\par\vspace{0.5em}
    {\fontsize{11}{13}\selectfont\itshape Communicated and edited by Boming Jia\par}}
\predate{}
\postdate{}
\renewcommand{\maketitlehookb}{\vspace{2.25em}}
\renewcommand{\maketitlehookd}{\vspace{-0.75em}}

\begin{document}
\maketitle

\begin{abstract}
    Let $G$ be a simply connected simple complex algebraic group, and let $U$ be a maximal unipotent subgroup. Assume that $G$ is not of type $A_1$ or $A_2$. We construct a $G$-equivariant closed Poisson embedding of the minimal nilpotent orbit closure into $T^*(G/U)^{\mathrm{aff}}$. The image of the minimal orbit is a symplectic leaf contained in the closure of every positive-dimensional symplectic leaf. It is the unique positive-dimensional leaf of smallest dimension.
\end{abstract}

\section{Introduction}

Let $G$ be a simply connected simple complex algebraic group with Lie algebra $\g$. Fix a Borel subgroup $B=TU$ with maximal torus $T$ and unipotent radical $U$. The affinization $T^*(G/U)^{\mathrm{aff}}=\Spec\C[T^*(G/U)]$ is a Poisson variety of finite type \cite[Lemma~3.6.2]{GinzburgRiche}.

We identify $\g$ with $\g^*$ by the Killing form and equip $\g^*$ with the Kirillov--Kostant--Souriau (KKS) Poisson structure. Let $\theta$ be the highest root, and choose $0\ne e_\theta\in\g_\theta$. The minimal nonzero nilpotent orbit $\Omin=Ge_\theta$ is a symplectic leaf of smallest positive dimension in $\g^*$.

\begin{theorem}\label{thm:main}
    Assume that $G$ is not of type $A_1$ or $A_2$. There is a $G$-equivariant closed Poisson embedding
    $$
        \overline{\Omin}\overset{j}{\hookrightarrow} T^*(G/U)^{\mathrm{aff}}
    $$
    such that $j(\Omin)$ is the unique positive-dimensional symplectic leaf of smallest dimension.
\end{theorem}

We will also show that $j(\Omin)$ is contained in the closure of every positive-dimensional symplectic leaf.

Both exclusions are necessary. In type $A_1$, we have $G/U\simeq\C^2\setminus\{0\}$ and $T^*(G/U)^{\mathrm{aff}}\simeq T^*\C^2$ with its standard symplectic form. In type $A_2$, it is Poisson isomorphic to the minimal nilpotent orbit closure in $\mathfrak{so}_8$ \cite[Corollary~3.8 and Theorem~6.2]{JiaPublished}. The minimal positive leaf dimensions are $4$ and $10$, while $\dim\Omin$ is $2$ and $4$, respectively.

\medskip
\noindent\textbf{Human Editor's Remark.}
This paper was generated by GPT-6 Astra and revised together with the Human Editor (HE), Boming Jia.
HE has found no mathematical errors in this paper.
If any errors are found later, HE will take full responsibility.
If the proof is indeed correct, HE does not want any credit for HE's negligible mathematical contribution.
\par\vspace{0.5em}
\noindent This work is licensed under a Creative Commons Attribution 4.0
International License (\mbox{\href{https://creativecommons.org/licenses/by/4.0/}{CC BY 4.0}}).

\setlength{\abovedisplayskip}{7pt plus 2pt minus 2pt}
\setlength{\belowdisplayskip}{7pt plus 2pt minus 2pt}
\setlength{\abovedisplayshortskip}{0pt plus 2pt}
\setlength{\belowdisplayshortskip}{5pt plus 2pt minus 2pt}
\AddToHook{env/theorem/before}{\Needspace{25\baselineskip}}
\newpage
\vspace{2em}
\section{Functions on the cotangent bundle}\label{sec:functions}

Put $\tfrak=\operatorname{Lie}(T)$, $\mathfrak b=\operatorname{Lie}(B)$, and $\mathfrak u=\operatorname{Lie}(U)$. The Killing form $\kappa$ identifies $\g$ with $\g^*$ and $\mathfrak b/\mathfrak u$ with $\tfrak^*$. Since $\mathfrak u^\perp=\mathfrak b$, we have
$$
    T^*(G/U)=G\times_U\mathfrak b,\qquad
    [g,\xi]=[gu,\Ad(u^{-1})\xi]\quad(u\in U).
$$
Our cotangent Poisson convention is
$$
    \{q_i,p_j\}=\delta_{ij},\qquad
    \{q_i,q_j\}=\{p_i,p_j\}=0.
$$
On $\g^*$, the KKS bracket is
$$
    \{a,b\}=[a,b]\qquad(a,b\in\g),
$$
where $a$ and $b$ are viewed as linear functions.

The cotangent lifts of the left $G$- and right $T$-actions are
$$
    g_0\cdot[g,\xi]=[g_0g,\xi],\qquad
    [g,\xi]s=[gs,\Ad(s^{-1})\xi].
$$
These commuting Poisson actions extend to the affinization. On regular functions, we use inverse pullback on the left and direct pullback on the right. Thus we have
$$
    (g_0\cdot f)(x)=f(g_0^{-1}\cdot x),\qquad
    (f\cdot s)(x)=f(xs).
$$
Throughout the paper, $T$-invariants and $T$-fixed points refer to the right action.

The moment maps for these actions are
$$
    \mu([g,\xi])=\Ad(g)\xi\in\g^*,\qquad
    \nu([g,\xi])=\xi\bmod\mathfrak u\in\tfrak^*.
$$
Their coordinate functions extend to the affinization. The map $\mu$ is $G$-equivariant and $T$-invariant, and $\nu$ is invariant under both actions.

For $a\in\g$ and $h\in\tfrak$, put
$$
    \mu_a=\langle\mu,a\rangle,\qquad
    \nu_h=\langle\nu,h\rangle.
$$
The left $\g$-action and right $\tfrak$-action on functions are
\begin{align*}
    (a\cdot f)(x)
    &=\left.\frac{d}{dt}\right|_{t=0}f(\exp(-ta)\cdot x),\\
    (f\cdot h)(x)
    &=\left.\frac{d}{dt}\right|_{t=0}f(x\exp(th)).
\end{align*}

\begin{proposition}\label{prop:infinitesimal-actions}\label{prop:moment-brackets}
    For $a,b\in\g$, $h,h'\in\tfrak$, and $f\in\C[T^*(G/U)]$, we have
    \begin{equation}\label{eq:infinitesimal-action}
        a\cdot f=\{\mu_a,f\},\qquad
        f\cdot h=\{f,\nu_h\},
    \end{equation}
    and
    \begin{equation}\label{eq:moment-brackets}
        \{\mu_a,\mu_b\}=\mu_{[a,b]},\qquad
        \{\mu_a,\nu_h\}=0,\qquad
        \{\nu_h,\nu_{h'}\}=0.
    \end{equation}
\end{proposition}

\begin{proof}
    For a vector field $v$ on $G/U$, let $\widehat v$ be its cotangent lift and $p(v)$ its fiber-linear function. In cotangent coordinates, write
    $$
        v=\sum_i v_i(q)\partial_{q_i},\qquad
        p(v)=\sum_i p_iv_i(q).
    $$
    Then we have
    $$
        \{p(v),q_i\}=-v_i(q),\qquad
        \{p(v),p_i\}=\sum_j p_j\frac{\partial v_j}{\partial q_i}.
    $$
    Thus $\{p(v),-\}=-\widehat v$. For the left vector field defined by $a$, the fiber pairing at $[g,\xi]$ is
    $$
        \kappa(\xi,\Ad(g^{-1})a)=\mu_a([g,\xi]).
    $$
    For the right vector field defined by $h$, it is
    $$
        \kappa(\xi,h)=\nu_h([g,\xi]).
    $$
    By our pullback conventions, we have \eqref{eq:infinitesimal-action}.

    By the equivariance of $\mu$, we have
    $$
        g\cdot\mu_b=\mu_{\Ad(g)b}.
    $$
    Differentiating, we have
    $$
        a\cdot\mu_b=\mu_{[a,b]}.
    $$
    Since $\nu$ is $G$- and $T$-invariant, the remaining brackets vanish.
\end{proof}

\begin{proposition}\label{prop:G-invariants}
    We have
    $$
        \C[T^*(G/U)]^G=\C[\nu_h\mid h\in\tfrak].
    $$
\end{proposition}

\begin{proof}
    Restriction to $[1,\xi]$ identifies $\C[T^*(G/U)]^G$ with $\C[\mathfrak b]^U$. By \cite[Section~7.3, before Lemma~7.3.3]{GinzburgKazhdan}, we have an isomorphism induced by the projection $\mathfrak b\to\mathfrak b/\mathfrak u$,
    $$
        \C[\mathfrak b/\mathfrak u]\simeq\C[\mathfrak b]^U.
    $$
    Its linear functions pull back to the $\nu_h$.
\end{proof}

Write $X^*(T)$ and $X_*(T)$ for the character and cocharacter lattices, with their natural pairing $\langle-,-\rangle$. For $\lambda\in X^*(T)$, write $d\lambda\in\tfrak^*$ for its differential at the identity.

\begin{definition}\label{def:weights}
    A function $f$ has \emph{right weight} $\lambda$ if
    $$
        f\cdot s=\lambda(s)f\qquad(s\in T).
    $$
    Denote this weight space by $\C[T^*(G/U)]_\lambda$. The function $f$ has \emph{left weight} $\chi$ if
    $$
        s\cdot f=\chi(s)f\qquad(s\in T),
    $$
    and \emph{fiber degree} $k$ if
    $$
        f([g,t\xi])=t^kf([g,\xi])\qquad(t\in\C^\times).
    $$
\end{definition}

The same right-weight convention defines $\C[G/U]_\lambda$. Pullback identifies it with the fiber-degree-zero part of $\C[T^*(G/U)]_\lambda$. For a function of left weight $\chi$ and right weight $\lambda$, we have
$$
    h\cdot f=d\chi(h)f,\qquad
    f\cdot h=d\lambda(h)f,\qquad
    \{\nu_h,f\}=-d\lambda(h)f.
$$

Let $\Phi$ be the root system of $(G,T)$, with positive roots $\Phi^+$ determined by $B$ and simple roots $\alpha_1,\ldots,\alpha_r$. We use the root-space notation $\g_\alpha$ and write $\alpha^\vee\in X_*(T)$ for the coroot of $\alpha$. Put
$$
    2\rho^\vee=\sum_{\alpha\in\Phi^+}\alpha^\vee,
$$
and let $W$ be the Weyl group. Dominance is taken with respect to $\Phi^+$.

\begin{definition}\label{def:conical-action}
    For $t\in\C^\times$, set
    $$
        c_t[g,\xi]
        =\bigl[g(2\rho^\vee)(t),\,
        t^2\Ad((2\rho^\vee)(t)^{-1})\xi\bigr].
    $$
    A nonzero regular function $f$ has \emph{$c$-degree} $d$ if
    $$
        c_t^*f=f\circ c_t=t^df\qquad(t\in\C^\times).
    $$
    We write $\deg_c f=d$.
\end{definition}

This action is the product of fiber scaling by $t^2$ and the right action of $(2\rho^\vee)(t)$. It commutes with $G\times T$ and extends to the affinization. In type $A$, this is the action in \cite[Definition~5.3 and Proposition~5.5]{JiaVOne}.

\begin{proposition}\label{prop:c-degree}\label{prop:c-bracket}
    For a nonzero function $f$ of fiber degree $k$ and right weight $\lambda$, we have
    \begin{equation}\label{eq:c-degree}
        \deg_c f=2k+\langle\lambda,2\rho^\vee\rangle.
    \end{equation}
    The Poisson bracket has $c$-degree $-2$.
\end{proposition}

\begin{proof}
    The two factors defining $c_t$ act on $f$ by $t^{2k}$ and $t^{\langle\lambda,2\rho^\vee\rangle}$. Thus we have \eqref{eq:c-degree}. By fiber scaling and the Poisson property of the right $T$-action, we have
    $$
        \{c_t^*f,c_t^*g\}=t^2c_t^*\{f,g\}.
    $$
    Thus the bracket has degree $-2$.
\end{proof}

\Needspace{8\baselineskip}
The following two results describe the right weight spaces as $G$-modules.

\begin{proposition}[Ginzburg--Riche]\label{prop:weight-spaces}
    For every dominant weight $\lambda$, we have
    $$
        \C[T^*(G/U)]_\lambda
        =\C[\mu_a,\nu_h\mid a\in\g,\ h\in\tfrak]\,\C[G/U]_\lambda.
    $$
    In particular, we have
    $$
        \C[T^*(G/U)]^T=\C[\mu_a,\nu_h].
    $$
\end{proposition}

\begin{proof}
    By the proof of \cite[Lemma~3.6.2]{GinzburgRiche}, we have a surjective multiplication map
    $$
        \Sym(\g\oplus\tfrak)\otimes\C[G/U]_\lambda
        \longrightarrow\C[T^*(G/U)]_\lambda,
    $$
    where the symmetric algebra acts through $\mu$ and $\nu$. This proves the first assertion. The second assertion is the case $\lambda=0$, since $\C[G/U]_0=\C[G/B]=\C$.
\end{proof}

\begin{proposition}[Ginzburg--Riche]\label{prop:weyl-symmetry}
    For $\lambda\in X^*(T)$ and $w\in W$, the spaces $\C[T^*(G/U)]_\lambda$ and $\C[T^*(G/U)]_{w\lambda}$ are isomorphic as $c$-graded $G$-modules.
\end{proposition}

\begin{proof}
    The projection $G\times_U\mathfrak b\to G\times_B\mathfrak b$ is a right $T$-torsor. Let $\mathcal O(\lambda)$ be the pullback to $G\times_B\mathfrak b$ of the line bundle on $G/B$ associated with $-\lambda$. Functions of right weight $\lambda$ correspond to its sections. Thus we have
    $$
        \C[T^*(G/U)]_\lambda
        \simeq\Gamma(G\times_B\mathfrak b,\mathcal O(\lambda)).
    $$
    These are the section spaces in \cite[(5.2.3)]{GinzburgRiche}. After forgetting the $\Sym(\tfrak)$-module structures, we have a $G$-equivariant isomorphism preserving the shifted grading.

    In \cite[Sections~1.10 and~2.4]{GinzburgRiche}, the action scales $\mathfrak b$ by $t^{-2}$ and acts on sections by inverse pullback. Thus fiber degree $k$ corresponds to degree $2k$. In weight $\lambda$, the grading is shifted by $\langle\lambda,2\rho^\vee\rangle$. If the image has fiber degree $k'$, we have
    $$
        2k+\langle\lambda,2\rho^\vee\rangle
        =2k'+\langle w\lambda,2\rho^\vee\rangle.
    $$
    By \eqref{eq:c-degree}, the two sides are their $c$-degrees. Since the isomorphism is $G$-equivariant, it also preserves every left $T$-weight.
\end{proof}


\vspace{2em}
\section{Symplectic leaves}\label{sec:leaves}

We recall the facts about Poisson subvarieties and leaf decompositions used below, following \cite{BrownGordon}.

Let $A$ be a finitely generated commutative $\C$-algebra. All points in this section are closed points.

\begin{definition}\label{def:poisson}
    A \emph{Poisson bracket} on $A$ is a $\C$-bilinear Lie bracket satisfying
    $$
        \{f,gh\}=\{f,g\}h+g\{f,h\}\qquad(f,g,h\in A).
    $$
    For $f\in A$, the derivation $H_f=\{f,-\}$ is its \emph{Hamiltonian vector field}. An ideal $I\subset A$ is \emph{Poisson} if $\{I,A\}\subset I$. A closed subvariety of $\Spec A$ is Poisson if its vanishing ideal is Poisson.
\end{definition}

Assume from now on that $A$ has a Poisson bracket.

\begin{definition}\label{def:poisson-rank}
    For $z\in\Spec A$, let $\mathfrak m_z$ be its maximal ideal and write
    $$
        d_zf=[f-f(z)]\in\mathfrak m_z/\mathfrak m_z^2.
    $$
    The \emph{Poisson rank} at $z$ is the rank of the alternating form
    $$
        \mathfrak m_z/\mathfrak m_z^2\times
        \mathfrak m_z/\mathfrak m_z^2\xrightarrow{\pi_z}\C,
        \qquad \pi_z(d_zf,d_zg)=\{f,g\}(z).
    $$
\end{definition}

By the Leibniz rule, we have
$$
    \{\mathfrak m_z^2,A\}\subset\mathfrak m_z.
$$
Hence $\pi_z$ is well-defined. Its associated map
$$
    \mathfrak m_z/\mathfrak m_z^2\xrightarrow{\pi_z^\sharp}T_z(\Spec A)
$$
sends $d_zf$ to $H_f(z)$. If $f_1,\ldots,f_N$ generate $A$, their differentials define a surjection $\C^N\to\mathfrak m_z/\mathfrak m_z^2$. By pulling back $\pi_z$ along this surjection, we have
$$
    \operatorname{rank}\pi_z
    =\operatorname{rank}\bigl(\{f_i,f_j\}(z)\bigr)_{1\le i,j\le N},
$$
as in \cite[Section~3.1]{BrownGordon}.

\begin{proposition}[{\cite[Sections~2.4 and~3.5]{BrownGordon}}]\label{prop:poisson-subvarieties}
    The bracket extends uniquely to every localization of $A$ and induces a Poisson bracket on $A/I$ for every Poisson ideal $I$. The radical $\sqrt I$ is Poisson. The reduced singular locus of $\Spec(A/\sqrt I)$ is a Poisson subvariety of $\Spec A$.
\end{proposition}

\begin{proof}
    Localization is defined by
    $$
        \{s^{-1},f\}=-s^{-2}\{s,f\},
    $$
    and the quotient bracket is
    $$
        \{f+I,g+I\}=\{f,g\}+I.
    $$

    Pass to $A/I$ and let $D$ be a Hamiltonian derivation. If $f^n=0$, then
    $$
        0=D^n(f^n)\equiv n!(Df)^n\pmod{(f)}.
    $$
    Every other term contains an undifferentiated factor $f$. Hence $(Df)^n\in(f)$ and $(Df)^{n^2}=0$. Thus the nilradical is preserved.

    The derivation $D$ descends to $A/\sqrt I$. Embed $\Spec(A/\sqrt I)$ in affine space and lift $D$ to a polynomial vector field preserving its defining ideal. For generators $q_i$ of that ideal, write
    $$
        Dq_i=\sum_j a_{ij}q_j.
    $$
    Along the local holomorphic flow, these functions satisfy a homogeneous linear differential system. Thus the flow preserves the variety. Since the flow is locally invertible, it also preserves the singular locus. By differentiation, $D$ preserves the vanishing ideal of that locus, as in \cite[Section~3.5]{BrownGordon}.
\end{proof}

\begin{definition}\label{def:symplectic-leaf}
    Set
    $$
        X_0=(\Spec A)_{\mathrm{red}},\qquad
        X_{i+1}=(\operatorname{Sing}X_i)_{\mathrm{red}},\qquad
        S_i=X_i\setminus X_{i+1}.
    $$
    The $X_i$ are Poisson by Proposition~\ref{prop:poisson-subvarieties}. Their dimensions strictly decrease, so the nonempty $S_i$ form a finite decomposition into smooth locally closed Poisson varieties. Following \cite[Section~3.5]{BrownGordon}, a \emph{symplectic leaf} of $\Spec A$ is a maximal connected immersed complex submanifold $L$ of one of the $S_i$ satisfying
    $$
        T_zL=\operatorname{span}_{\C}
        \{H_f(z)\mid f\in\mathcal O^{\mathrm{an}}_{S_i,z}\}
        \qquad(z\in L).
    $$
    Here the bracket is extended to germs of holomorphic functions on $S_i$.
\end{definition}

\begin{proposition}[{\cite[Section~3.5]{BrownGordon}}]\label{prop:leaf-tangent}
    Each point $z$ lies in a unique symplectic leaf $L$. Its tangent space and dimension satisfy
    \begin{equation}\label{eq:leaf-rank}
        T_zL=\operatorname{span}_{\C}\{H_f(z)\mid f\in A\}
        =\operatorname{im}\pi_z^\sharp,\qquad
        \dim L=\operatorname{rank}\pi_z.
    \end{equation}
    The induced bracket on $L$ is nondegenerate and satisfies
    $$
        \{f,g\}|_L=\{f|_L,g|_L\}_L\qquad(f,g\in A).
    $$
\end{proposition}

\begin{proof}
    Existence and uniqueness on each $S_i$ follow from the Poisson foliation theorem, as recalled in \cite[Section~3.5]{BrownGordon}. Let $z\in S_i$, and let $I_i\subset A$ define $X_i$. Since $S_i$ is open in $X_i$,
    $$
        T_z^*S_i=\mathfrak m_z/(\mathfrak m_z^2+I_i).
    $$
    The ideal $I_i$ is Poisson, so its image in $\mathfrak m_z/\mathfrak m_z^2$ lies in the radical of $\pi_z$. Quotienting by this subspace does not change the rank. The differentials of algebra generators span this cotangent space, so their Hamiltonian fields span $T_zL$. This proves \eqref{eq:leaf-rank}.

    By the holomorphic inverse function theorem, compositions of flows of a local Hamiltonian frame parametrize $L$ near each point. Thus each flow orbit is open in $L$. Since $L$ is connected, it is a single flow orbit. The formula
    $$
        \omega_L(H_f,H_g)=\{f,g\}|_L
    $$
    is well-defined, since $H_f(z)=0$ implies $\{f,g\}(z)=0$ for every $g$. In a local Hamiltonian frame it defines a holomorphic two-form. A tangent vector pairing to zero with every $H_g$ annihilates every $g\in A$ and is therefore zero. Thus the form is nondegenerate. By the Jacobi identity, it is closed. Since
    $$
        \iota_{H_f}\omega_L=-d(f|_L),
    $$
    its Poisson bracket is the restriction of the ambient bracket.
\end{proof}

In particular, if $f|_L$ is constant, then $\{f,g\}|_L=0$ for every $g\in A$, so $H_f$ vanishes on $L$.

\begin{proposition}[{\cite[Lemma~3.5 and its proof]{BrownGordon}}]\label{prop:leaf-closure}
    The Zariski closure of a leaf $L$ is Poisson. For a closed Poisson subvariety $Y\subset\Spec A$, the leaves of $Y$ are precisely the ambient leaves contained in $Y$.
\end{proposition}

\begin{proof}
    If $f$ vanishes on $\overline L$, then
    $$
        \{g,f\}|_L=\{g|_L,0\}_L=0\qquad(g\in A).
    $$
    Thus the vanishing ideal of $\overline L$ is Poisson.

    Let $I$ define $Y$. Each $H_f$ descends to $H_{f+I}$ on $Y$. By the proof of Proposition~\ref{prop:poisson-subvarieties}, its local flow preserves $Y$. The algebra generators of $A$ restrict to generators of $A/I$. By Proposition~\ref{prop:leaf-tangent}, their flows generate the leaves in both varieties. Since these flows agree on $Y$, the leaves agree.
\end{proof}

\begin{proposition}[{\cite[Proposition~3.7]{BrownGordon}}]\label{prop:finite-leaves}
    If $\Spec A$ has finitely many symplectic leaves, each leaf is a smooth locally closed algebraic subvariety. The leaves are precisely the connected components, equivalently the irreducible components, of the loci of constant Poisson rank.
\end{proposition}

\begin{proof}
    Passing to $A_{\mathrm{red}}$ changes neither the leaves nor their ranks, since the nilradical is Poisson. By \cite[Proposition~3.7(1), (2)(a),(c)]{BrownGordon}, the leaves are the irreducible components of the sets
    $$
        \{z\in\Spec A\mid\operatorname{rank}\pi_z=r\},
        \qquad r\in\mathbb Z_{\ge0}.
    $$
    These sets are locally closed, since rank at most $r$ is defined by the $(r+1)\times(r+1)$ minors of the bracket matrix. Their irreducible components are the leaves, so they are disjoint and finite in number. Each is therefore open and closed in its rank locus, hence also a connected component.
\end{proof}

\begin{proposition}\label{prop:implosion-finiteness}
    The variety $T^*(G/U)^{\mathrm{aff}}$ has finitely many symplectic leaves.
\end{proposition}

\begin{proof}
    By \cite[Theorem~1.1]{Gannon}, the variety $T^*(G/U)^{\mathrm{aff}}$ has symplectic singularities. Such a variety has finitely many symplectic leaves by Kaledin's theorem \cite[Theorem~2.3]{Kaledin}.
\end{proof}

\Needspace{8\baselineskip}
\begin{proposition}\label{prop:minimal-adjoint-dimension}
    The orbit $\Omin$ has the smallest dimension among all nonzero adjoint orbits in $\g$. Equivalently, it is a symplectic leaf of smallest positive dimension in $\g^*$.
\end{proposition}

\begin{proof}
    Let $0\ne x\in\g$. By the Borel fixed-point theorem, the projective orbit closure $\overline{G\cdot[x]}\subset\mathbb P(\g)$ contains a $B$-fixed point. Since the adjoint $G$-module $\g$ is simple, its only $B$-stable line is $\C e_\theta$. Thus we have
    $$
        G\cdot[e_\theta]\subset\overline{G\cdot[x]}.
    $$
    The orbit $\Omin$ is stable under nonzero scalars, so
    $$
        \dim\Omin-1
        =\dim(G\cdot[e_\theta])
        \le\dim(G\cdot[x])
        \le\dim(Gx).
    $$
    Both $\dim\Omin$ and $\dim(Gx)$ are even, since adjoint orbits carry the KKS symplectic form. Hence we have
    \[
        \dim\Omin\le\dim(Gx).\qedhere
    \]
\end{proof}

\begin{remark}\label{rem:type-a-minimal-dimension}
    A symplectic leaf of smallest positive dimension in $\g^*$ need not be unique. In type $A_{n-1}$ with $n\ge2$, the orbit $\Omin\subset\mathfrak{sl}_n$ consists of rank-one nilpotent matrices and has dimension $2n-2$. The semisimple element
    $$
        s=\operatorname{diag}(n-1,-1,\ldots,-1)
    $$
    has a centralizer in $\mathfrak{sl}_n$ of dimension $(n-1)^2$, so
    $$
        \dim(Gs)=(n^2-1)-(n-1)^2=2n-2=\dim\Omin.
    $$
    Thus $Gs$ is another symplectic leaf of the same minimal positive dimension.
\end{remark}

\vspace{2em}
\section{The minimal leaf}\label{sec:minimal-leaf}

Throughout this section, assume that $G$ is not of type $A_1$ or $A_2$.

Recall that
$$
    \overline{\Omin}=\Omin\sqcup\{0\}
$$
by \cite[Section~4.3]{CollingwoodMcGovern}. Since
$$
    \Ad(\theta^\vee(t))e_\theta=t^2e_\theta,
$$
every nonzero multiple of $e_\theta$ belongs to $\Omin$. Choose $e_{-\theta}\in\g_{-\theta}$ with $\kappa(e_\theta,e_{-\theta})=1$.

For $t\in\C^\times$, define
$$
    \gamma_t x=\theta^\vee(t^{-1})\cdot c_t x.
$$
Since $c$ commutes with the left $G$-action, this defines an algebraic $\C^\times$-action. It commutes with the right $T$-action, so each right weight space is $\gamma$-graded. By the equivariance of $\mu$ and the invariance of $\nu$, we have
$$
    \mu(\gamma_t x)=t^2\Ad(\theta^\vee(t^{-1}))\mu(x),\qquad
    \nu(\gamma_t x)=t^2\nu(x).
$$
Write $\deg_\gamma f=m$ if $f\ne0$ and
$$
    \gamma_t^*f=f\circ\gamma_t=t^mf\qquad(t\in\C^\times).
$$
If $f$ has $c$-degree $d$ and left weight $\chi$, then
\begin{equation}\label{eq:gamma-degree}
    \deg_\gamma f=d+\langle\chi,\theta^\vee\rangle.
\end{equation}
By the left pullback convention, we have
$$
    \gamma_t^*f=c_t^*(\theta^\vee(t)\cdot f).
$$
The left $G$-action is Poisson. By Proposition~\ref{prop:c-bracket}, we also have
\begin{equation}\label{eq:scaled-bracket}
    \{\gamma_t^*f,\gamma_t^*g\}
    =t^2\gamma_t^*\{f,g\}\qquad(t\ne0).
\end{equation}
Thus the bracket has $\gamma$-degree $-2$.

\begin{lemma}\label{lem:root-estimates}
    The coroot $\theta^\vee$ is dominant, and we have
    $$
        \langle\beta,\theta^\vee\rangle\ge-2\qquad(\beta\in\Phi),
    $$
    with equality if and only if $\beta=-\theta$. For every nonzero dominant weight $\lambda$, we have
    $$
        \langle\lambda,2\rho^\vee-\theta^\vee\rangle\ge2.
    $$
\end{lemma}

\begin{proof}
    Fix a $W$-invariant positive definite inner product $(\, ,\, )$ on the real span of $\Phi$. By the root-string property and the maximality of $\theta$, we have
    $$
        (\theta,\alpha_i)\ge0\qquad(1\le i\le r).
    $$
    Thus $\theta^\vee$ is dominant. Since $\theta$ is long \cite[Chapter~VI, Section~1]{Bourbaki}, we have
    $$
        \bigl|\langle\beta,\theta^\vee\rangle\bigr|
        =\frac{2|(\beta,\theta)|}{(\theta,\theta)}
        \le\frac{2\sqrt{(\beta,\beta)}}{\sqrt{(\theta,\theta)}}\le2.
    $$
    Equality with value $-2$ holds only if $\beta=-\theta$.

    Since the Dynkin diagram is connected and $r\ge2$, we can choose adjacent simple roots $\alpha_i,\alpha_j$ with
    $$
        \langle\lambda,\alpha_i^\vee\rangle>0.
    $$
    By the root-string property, we have
    $$
        \alpha_i+\alpha_j\in\Phi^+.
    $$
    Both $\alpha_i$ and $\alpha_i+\alpha_j$ differ from $\theta$. For $r\ge3$, this follows from the full support of $\theta$. Indeed, an omitted simple root adjacent to its support would have negative inner product with $\theta$, contradicting maximality by the root-string property. For $r=2$, the exclusion of $A_2$ means that a Cartan integer is at most $-2$, so a root string has a root of height at least three. Hence $\theta$ has height at least three.

    The coroot $(\alpha_i+\alpha_j)^\vee$ is a positive integral combination of $\alpha_i^\vee$ and $\alpha_j^\vee$. By dominance and integrality of $\lambda$, both $\alpha_i$ and $\alpha_i+\alpha_j$ contribute at least $1$ to
    $$
        \langle\lambda,2\rho^\vee-\theta^\vee\rangle
        =\sum_{\beta\in\Phi^+\setminus\{\theta\}}
        \langle\lambda,\beta^\vee\rangle.
    $$
    All other summands are nonnegative.
\end{proof}

\setlength{\abovedisplayskip}{12pt plus 2pt minus 2pt}
\setlength{\belowdisplayskip}{12pt plus 2pt minus 2pt}
\setlength{\abovedisplayshortskip}{0pt plus 2pt}
\setlength{\belowdisplayshortskip}{10pt plus 2pt minus 2pt}
\begin{theorem}\label{thm:gamma-grading}
    The $\gamma$-grading on $\C[T^*(G/U)]$ is nonnegative. Its degree-zero algebra is $\C[\mu_{e_{-\theta}}]$. Every right weight space of nonzero weight is concentrated in $\gamma$-degrees at least $2$.
\end{theorem}

\begin{proof}
    By Proposition~\ref{prop:weight-spaces}, the right weight-zero algebra is generated by the moment functions. They all have $c$-degree $2$. Let $a\in\g_\beta$ be a root vector. By equivariance, we have $g\cdot\mu_a=\mu_{\Ad(g)a}$. By \eqref{eq:gamma-degree}, we have
    $$
        \deg_\gamma\mu_a=2+\langle\beta,\theta^\vee\rangle.
    $$
    The functions $\mu_h$ and $\nu_h$, with $h\in\tfrak$, have $\gamma$-degree $2$. By Lemma~\ref{lem:root-estimates}, all these degrees are nonnegative, and only $\g_{-\theta}$ contributes degree-zero moment functions. A degree-zero monomial can involve only degree-zero generators. Thus the degree-zero part of the right invariant algebra is $\C[\mu_{e_{-\theta}}]$. This is a polynomial algebra since
    $$
        \mu_{e_{-\theta}}([1,se_\theta])=s\qquad(s\in\C).
    $$

    Let $\lambda\ne0$ be dominant. By Borel--Weil, $\C[G/U]_\lambda$ is the dual of the irreducible module $V_\lambda$ of highest weight $\lambda$. For a highest-weight vector $v_\lambda$ and $\ell\in V_\lambda^*$, the function $gU\mapsto\ell(gv_\lambda)$ has right weight $\lambda$. The left action on $\ell$ is the dual representation. Thus a nonzero left weight vector $f\in\C[G/U]_\lambda$ has weight
    $$
        \chi=-\lambda+\sum_i n_i\alpha_i,\qquad n_i\ge0.
    $$
    Since $f$ has fiber degree zero, we have
    \begin{align*}
        \deg_\gamma f
        &=\langle\lambda,2\rho^\vee-\theta^\vee\rangle
          +\sum_i n_i\langle\alpha_i,\theta^\vee\rangle\\
        &\ge\langle\lambda,2\rho^\vee-\theta^\vee\rangle\ge2.
    \end{align*}
    By Proposition~\ref{prop:weight-spaces}, the space $\C[T^*(G/U)]_\lambda$ is generated by $\C[G/U]_\lambda$ over the right invariant algebra. Since that algebra has nonnegative degrees, the same bound holds on $\C[T^*(G/U)]_\lambda$.

    For arbitrary $\lambda\ne0$, choose $w\in W$ with $w\lambda$ dominant. By Proposition~\ref{prop:weyl-symmetry}, we have an isomorphism preserving $c$-degree and left $T$-weights. By \eqref{eq:gamma-degree}, it also preserves $\gamma$-degree. Hence the same bound holds for every nonzero right weight.
\end{proof}

\begin{corollary}\label{cor:contraction}
    The action $\gamma$ extends to a unique morphism
    $$
        \mathbb A^1\times T^*(G/U)^{\mathrm{aff}}
        \longrightarrow T^*(G/U)^{\mathrm{aff}}.
    $$
    The fixed point subscheme of $\gamma$ is an affine line with coordinate $\mu_{e_{-\theta}}$, and $\gamma_0$ is a retraction onto this line.
\end{corollary}

\begin{proof}
    By Theorem~\ref{thm:gamma-grading}, the rule
    $$
        f\longmapsto t^m\otimes f
    $$
    in degree $m$ defines an algebra homomorphism and hence the extension. Multiplication by $t$ in $\C[t]\otimes\C[T^*(G/U)]$ is injective. Thus two extensions that agree after inverting $t$ already agree.

    For a homogeneous function $f$ of degree $m$, the fixed point equations are
    $$
        (t^m-1)f=0\qquad(t\in\C^\times).
    $$
    Thus the fixed point ideal is the positive-degree ideal, and its quotient is $\C[\mu_{e_{-\theta}}]$. The map $\gamma_0^*$ is the degree-zero projection. It is the identity on this quotient, so $\gamma_0$ is a retraction.
\end{proof}

\begin{lemma}\label{lem:fixed-scheme-immersion}
    The restriction
    $$
        (\nu^{-1}(0))^T\overset{\mu}{\longrightarrow}\g^*
    $$
    is a $G$-equivariant closed immersion.
\end{lemma}

\begin{proof}
    The zero fiber and the fixed points are taken scheme-theoretically. Their defining ideal in $\C[T^*(G/U)]$ is
    $$
        I=\bigl\langle\nu_h,\ \C[T^*(G/U)]_\lambda
        \mid h\in\tfrak,\ \lambda\ne0\bigr\rangle.
    $$
    Since the $\nu_h$ are invariant, the coordinate ring of the zero fiber inherits the right weight decomposition. For a function of right weight $\lambda$, the fixed point equations are
    $$
        (\lambda(s)-1)f=0\qquad(s\in T).
    $$
    If $\lambda\ne0$, choose $s$ with $\lambda(s)\ne1$. Then $f=0$. Conversely, the action on the quotient by $I$ is trivial. This argument holds over every commutative $\C$-algebra.

    By Proposition~\ref{prop:weight-spaces}, the quotient by $I$ is generated by the images of the $\mu_a$. Thus $\mu^*$ induces a surjection
    $$
        \C[\g^*]\longrightarrow\C[T^*(G/U)]/I,
    $$
    so the restriction of $\mu$ is a closed immersion. The left $G$-action commutes with the right $T$-action and fixes each $\nu_h$. Hence $I$ is $G$-stable, and the restriction of $\mu$ is $G$-equivariant.
\end{proof}

\setlength{\abovedisplayskip}{7pt plus 2pt minus 2pt}
\setlength{\belowdisplayskip}{7pt plus 2pt minus 2pt}
\setlength{\abovedisplayshortskip}{0pt plus 2pt}
\setlength{\belowdisplayshortskip}{5pt plus 2pt minus 2pt}
\begin{lemma}\label{lem:implosion-leaves}
    Every symplectic leaf is $G$-stable and is preserved by $\gamma_t$ for $t\ne0$.
\end{lemma}

\begin{proof}
    The $G$-action is Poisson, so it preserves Poisson rank. For $t\ne0$, the pullback by $\gamma_t$ is an isomorphism on cotangent spaces. By \eqref{eq:scaled-bracket}, we have
    $$
        \pi_x\bigl(d_x(\gamma_t^*f),d_x(\gamma_t^*g)\bigr)
        =t^2\pi_{\gamma_t x}\bigl(d_{\gamma_t x}f,d_{\gamma_t x}g\bigr).
    $$
    Since $t^2\ne0$, we have
    $$
        \operatorname{rank}\pi_x=\operatorname{rank}\pi_{\gamma_t x}.
    $$
    This equality also holds at singular points.

    Both $G$ and $\C^\times$ are connected. By Propositions~\ref{prop:implosion-finiteness} and~\ref{prop:finite-leaves}, each leaf $L$ is a connected component of
    $$
        \{x\in T^*(G/U)^{\mathrm{aff}}\mid\operatorname{rank}\pi_x=\dim L\}.
    $$
    Each group orbit is connected and stays in that locus. Hence it stays in the component containing its initial point.
\end{proof}

\newpage
\setlength{\abovedisplayskip}{12pt plus 2pt minus 2pt}
\setlength{\belowdisplayskip}{12pt plus 2pt minus 2pt}
\setlength{\abovedisplayshortskip}{0pt plus 2pt}
\setlength{\belowdisplayshortskip}{10pt plus 2pt minus 2pt}
\begin{theorem}\label{thm:lift}
    There is a unique $G$-equivariant closed embedding
    $$
        \overline{\Omin}\overset{j}{\hookrightarrow}
        (\nu^{-1}(0))^T\subset T^*(G/U)^{\mathrm{aff}}
    $$
    such that $\mu\circ j$ is the inclusion into $\g^*$.
\end{theorem}

\Needspace{4\baselineskip}
\begin{proof}
    Each $\nu_h$ has $\gamma$-degree $2$, and each nonzero right weight space has degrees at least $2$. Thus $\gamma_0^*$ annihilates the generators of $I$ in Lemma~\ref{lem:fixed-scheme-immersion}. Since it is an algebra homomorphism, it annihilates $I$. Hence $\gamma_0$ factors scheme-theoretically through $(\nu^{-1}(0))^T$. By the degrees of the moment functions, we have
    \begin{equation}\label{eq:zero-projection}
        \gamma_0^*\mu_a=\kappa(e_\theta,a)\mu_{e_{-\theta}},
        \qquad \gamma_0^*\nu_h=0.
    \end{equation}
    Here the coefficient of $e_{-\theta}$ in $a$ is $\kappa(e_\theta,a)$, by our normalization. In particular, we have
    $$
        \mu(\gamma_0[1,e_\theta])=e_\theta.
    $$

    Let $J\subset\C[\g^*]$ be the ideal of the closed subscheme $\mu((\nu^{-1}(0))^T)$. This subscheme is $G$-stable and contains $e_\theta$. Hence it contains $\Omin$ on closed points. Every element of $J$ vanishes on $\Omin$ and its closure. Since $\overline{\Omin}$ is reduced, we have
    $$
        J\subset I(\overline{\Omin}).
    $$
    Since both ideals are $G$-stable, the quotient map
    $$
        \C[\g^*]/J\longrightarrow\C[\overline{\Omin}]
    $$
    defines a $G$-equivariant closed embedding $j$ with $\mu\circ j$ equal to the inclusion. Since the restriction of $\mu$ is a closed immersion, there is at most one morphism with this property.

    By \eqref{eq:zero-projection}, we have
    $$
        \mu(\gamma_0x)=\mu_{e_{-\theta}}(x)e_\theta.
    $$
    Both $\gamma_0$ and $x\mapsto j(\mu_{e_{-\theta}}(x)e_\theta)$ factor through $(\nu^{-1}(0))^T$. Their compositions with $\mu$ agree as morphisms to $\g^*$. Since the restriction of $\mu$ is a closed immersion, the two morphisms agree. Hence we have
    \begin{equation}\label{eq:limit}
        \gamma_0x=\lim_{t\to0}\gamma_t x
        =j\bigl(\mu_{e_{-\theta}}(x)e_\theta\bigr).
    \end{equation}

    On the fixed point line of Corollary~\ref{cor:contraction}, the function $\mu_{e_{-\theta}}$ is a coordinate. By \eqref{eq:zero-projection}, the restriction of $\mu$ is an isomorphism onto $\C e_\theta$. Thus the fixed point subscheme of $\gamma$ satisfies
    $$
        \bigl(T^*(G/U)^{\mathrm{aff}}\bigr)^\gamma=j(\C e_\theta).
    $$
    By the definition of $j$, we also have
    \begin{equation}\label{eq:embedding-functions}
        j^*\mu_a=\kappa(-,a),\qquad
        j^*\nu_h=0,\qquad j^*f=0
    \end{equation}
    for every $f$ of nonzero right weight.
\end{proof}

\newpage
\setlength{\abovedisplayskip}{12pt plus 2pt minus 2pt}
\setlength{\belowdisplayskip}{12pt plus 2pt minus 2pt}
\setlength{\abovedisplayshortskip}{0pt plus 2pt}
\setlength{\belowdisplayshortskip}{10pt plus 2pt minus 2pt}
\begin{theorem}\label{thm:minimal-orbit-leaf}
    The embedding $j$ is Poisson, and $j(\Omin)$ is a symplectic leaf of $T^*(G/U)^{\mathrm{aff}}$.
\end{theorem}

\begin{proof}
    By Proposition~\ref{prop:weight-spaces}, the coordinate algebra is generated by the moment functions and the functions of nonzero right weight. By the Leibniz rule, it suffices to verify the Poisson identity for $j^*$ on these generators. For two moment functions, the identity follows from \eqref{eq:moment-brackets}, \eqref{eq:embedding-functions}, and the KKS bracket. For a function $f$ of nonzero right weight, we have
    $$
        \{\mu_a,f\}=a\cdot f,\qquad
        \{\nu_h,f\}=-f\cdot h.
    $$
    Both brackets have the same right weight as $f$, so their pullbacks vanish. Since $j^*f=0$, the mixed Poisson identities hold.

    It remains to consider $f_1,f_2$ of nonzero right weights. Since $\gamma$ commutes with $T$, each $\gamma$-homogeneous component has the same right weight as the original function. By Theorem~\ref{thm:gamma-grading}, we have finite decompositions
    $$
        f_i=\sum_{m\ge2}f_{i,m}.
    $$
    Every nonzero bracket $\{f_{1,m},f_{2,n}\}$ has $\gamma$-degree $m+n-2\ge2$. By Corollary~\ref{cor:contraction}, we have
    $$
        \gamma_0^*\{f_1,f_2\}=0.
    $$
    Since $\gamma_0$ is a retraction onto $j(\C e_\theta)$, the bracket vanishes on this line.
    The $G$-action preserves right weights. Thus the preceding conclusion applies to $g^{-1}\cdot f_1$ and $g^{-1}\cdot f_2$ for every $g\in G$. Since this action is Poisson, we have
    $$
        \{f_1,f_2\}(g\cdot j(e_\theta))
        =\{g^{-1}\cdot f_1,g^{-1}\cdot f_2\}(j(e_\theta))=0.
    $$
    Thus $j^*\{f_1,f_2\}=0$ on $\overline{\Omin}$, since $\Omin$ is dense and its closure is reduced. This equals $\{j^*f_1,j^*f_2\}$ by \eqref{eq:embedding-functions}. Hence $j$ is Poisson.

    The image of $j$ is a closed Poisson subvariety isomorphic to $\overline{\Omin}$. Its open orbit is a KKS symplectic leaf. By Proposition~\ref{prop:leaf-closure}, $j(\Omin)$ is also an ambient symplectic leaf.
\end{proof}

\setlength{\abovedisplayskip}{5pt plus 2pt minus 2pt}
\setlength{\belowdisplayskip}{5pt plus 2pt minus 2pt}
\setlength{\abovedisplayshortskip}{0pt plus 2pt}
\setlength{\belowdisplayshortskip}{3pt plus 2pt minus 2pt}
\begin{lemma}\label{lem:moment-nonzero}
    No positive-dimensional symplectic leaf is contained in $\mu^{-1}(0)$.
\end{lemma}

\begin{proof}
    Suppose that $L\subset\mu^{-1}(0)$ is a symplectic leaf. The functions $\mu_a$ vanish on $L$. By Proposition~\ref{prop:leaf-tangent}, their Hamiltonian vector fields also vanish on $L$. By \eqref{eq:infinitesimal-action}, the stabilizer of each point has Lie algebra $\g$. Stabilizers are smooth over $\C$, so they have dimension $\dim G$. Since $G$ is connected, each stabilizer is $G$. Hence $G$ fixes $L$ pointwise.

    Every ambient regular function $f$ lies in a finite-dimensional $G$-submodule. By complete reducibility, write $f=f_0+f_1$, where $f_0$ is invariant and $f_1$ lies in the sum of the nontrivial simple summands. Evaluation at a $G$-fixed point is $G$-equivariant with trivial target, so it annihilates these summands. Hence $f|_L=f_0|_L$.

    Let $f,g$ be ambient regular functions. Choose invariant functions $f_0,g_0$ with the same restrictions to $L$. By Proposition~\ref{prop:G-invariants}, these are polynomials in the $\nu_h$, so $\{f_0,g_0\}=0$. By Proposition~\ref{prop:leaf-tangent}, we have
    $$
        \{f,g\}|_L
        =\{f|_L,g|_L\}_L
        =\{f_0|_L,g_0|_L\}_L=0.
    $$
    By \eqref{eq:leaf-rank}, we have $\dim L=0$.
\end{proof}

\newpage
\setlength{\abovedisplayskip}{5pt plus 2pt minus 2pt}
\setlength{\belowdisplayskip}{5pt plus 2pt minus 2pt}
\setlength{\abovedisplayshortskip}{0pt plus 2pt}
\setlength{\belowdisplayshortskip}{3pt plus 2pt minus 2pt}
\begin{theorem}\label{thm:leaf-closure}
    The closure of every positive-dimensional symplectic leaf contains $j(\overline{\Omin})$.
\end{theorem}

\begin{proof}
    Let $L$ be a positive-dimensional leaf. By Lemmas~\ref{lem:implosion-leaves} and~\ref{lem:moment-nonzero}, the linear span of $\mu(L)$ is a nonzero $G$-submodule of $\g$. Since $\g$ is simple, this span is $\g$. Choose $x\in L$ with
    $$
        \mu_{e_{-\theta}}(x)\ne0.
    $$

    The morphism $\mathbb A^1\to T^*(G/U)^{\mathrm{aff}}$, given by $t\mapsto\gamma_t x$, sends $\C^\times$ into $L$. The inverse image of $\overline L$ is closed and contains $\C^\times$, so it is all of $\mathbb A^1$. Thus $\gamma_0x\in\overline L$. By \eqref{eq:limit}, we have
    $$
        \gamma_0x=j\bigl(\mu_{e_{-\theta}}(x)e_\theta\bigr)\in j(\Omin).
    $$
    The closed set $\overline L$ is $G$-stable. It therefore contains $j(\Omin)$ and its closure $j(\overline{\Omin})$.
\end{proof}

\begin{proof}[Proof of Theorem~\ref{thm:main}]
    By Theorems~\ref{thm:lift} and~\ref{thm:minimal-orbit-leaf}, we have the closed Poisson embedding $j$ and the symplectic leaf $j(\Omin)$. Let $L$ be any other positive-dimensional leaf. Distinct leaves are disjoint. By Theorem~\ref{thm:leaf-closure}, we have
    $$
        j(\Omin)\subset\overline L\setminus L.
    $$
    By Proposition~\ref{prop:finite-leaves}, $L$ is irreducible and locally closed. Thus $L$ is open in its irreducible closure, and $\dim\overline L=\dim L$. The boundary is a proper closed subset of $\overline L$, so each of its irreducible components has smaller dimension. Since $j$ is a closed embedding, we have
    \[
        \dim\Omin\le\dim(\overline L\setminus L)<\dim L.\qedhere
    \]
\end{proof}

\begin{thebibliography}{Gan24}
    \setlength{\itemsep}{0pt}
    \setlength{\parsep}{0pt}
    \setlength{\parskip}{0pt}
    \interlinepenalty=10000

\bibitem[Bou02]{Bourbaki}
    N.~Bourbaki, \emph{Lie Groups and Lie Algebras, Chapters 4--6}, Elements of Mathematics, Springer-Verlag, Berlin, 2002.

\bibitem[BG03]{BrownGordon}
    K.~A.~Brown and I.~Gordon, Poisson orders, symplectic reflection algebras and representation theory, \emph{J. Reine Angew. Math.} \textbf{559} (2003), 193--216.

\bibitem[CM93]{CollingwoodMcGovern}
    D.~H.~Collingwood and W.~M.~McGovern, \emph{Nilpotent Orbits in Semisimple Lie Algebras}, Van Nostrand Reinhold, New York, 1993.

\bibitem[Gan24]{Gannon}
    T.~Gannon, Proof of the Ginzburg--Kazhdan conjecture, \emph{Adv. Math.} \textbf{448} (2024), Paper No.~109701, 28 pp.

\bibitem[GK22]{GinzburgKazhdan}
    V.~Ginzburg and D.~Kazhdan, Differential operators on $G/U$ and the Gelfand--Graev action, \emph{Adv. Math.} \textbf{403} (2022), Paper No.~108368, 48 pp.

\bibitem[GR15]{GinzburgRiche}
    V.~Ginzburg and S.~Riche, Differential operators on $G/U$ and the affine Grassmannian, \emph{J. Inst. Math. Jussieu} \textbf{14} (2015), no.~3, 493--575.

\bibitem[Jia21]{JiaVOne}
    B.~Jia, The geometry of the affine closure of $T^*(SL_n/U)$, \href{https://arxiv.org/abs/2112.08649v1}{arXiv 2112.08649v1} [math.RT] (2021).

\bibitem[Jia25]{JiaPublished}
    B.~Jia, The affine closure of $T^*(SL_n/U)$, \emph{J. Lie Theory} \textbf{35} (2025), no.~1, 83--100.

\bibitem[Kal06]{Kaledin}
    D.~Kaledin, Symplectic singularities from the Poisson point of view, \emph{J. Reine Angew. Math.} \textbf{600} (2006), 135--156.
\end{thebibliography}

\vfill
\noindent\hfill\begin{minipage}{350bp}
    \raggedleft\itshape
    ``Everybody understood that if the proof is correct\\
    then no other recognition is needed.''
    \par\vspace{0.25em}
    {\upshape --- Grigori Perelman}
\end{minipage}

\end{document}
